% !TEX program = lualatex
% Exemplo: tabela-verdade gerada automaticamente pelo pacote truthtable, que calcula os valores
% em Lua. Compilar com LuaLaTeX (no Overleaf: Menu > Compiler > LuaLaTeX).
% Compilar: latexmk -lualatex logica-truthtable.tex  Testado com: truthtable 0.1.0, TeX Live 2024 ou mais recente
%%BEGIN:truthtablepreambulo
\documentclass[12pt, a4paper]{article}
\usepackage[brazilian]{babel}   % LuaLaTeX: sem fontenc nem inputenc
\usepackage{amssymb}
\usepackage{truthtable}         % só funciona com LuaLaTeX
%%END:truthtablepreambulo

\begin{document}
\pagestyle{empty}

%%BEGIN:truthtableuso
% {variáveis}{cabeçalho delas}{fórmulas}{cabeçalho delas}{V}{F}
% & (e), | (ou), ! (não); >>(p;q) implica, <>(p;q) equivale,
% ^(p;q) ou exclusivo: argumentos separados por ponto e vírgula
\begin{center}
\begin{tabular}{cc|cccc}
  \truthtable{p,q}{$p$,$q$}
    {p & q, p | q, >>(p;q), <>(p;q)}
    {$p \land q$, $p \lor q$, $p \to q$, $p \leftrightarrow q$}
    {V}{F}
\end{tabular}
\end{center}
%%END:truthtableuso

%%BEGIN:truthtabletautologia
Com três variáveis, a tabela tem oito linhas. A última coluna,
só com V, mostra que $(p \to q) \land (q \to r) \to (p \to r)$
é uma tautologia:
\begin{center}
\begin{tabular}{ccc|ccc}
  \truthtable{p,q,r}{$p$,$q$,$r$}
    {>>(p;q), >>(q;r), >>(>>(p;q) & >>(q;r);>>(p;r))}
    {$p \to q$, $q \to r$,
     $(p \to q) \land (q \to r) \to (p \to r)$}
    {V}{F}
\end{tabular}
\end{center}
%%END:truthtabletautologia

\end{document}
